% ====================================================================
% Файл: tasks/01_mechanics/kim02.tex
% Тема: Задание №2 (Динамика)
% ====================================================================

% === ЗАДАЧА 1 ===
% Тег: #МЕХАНИКА #КИМ_02 #ВАРИАНТ_1
\begin{taskbasic}[code={\label{task:dem26_v1_n2}}]
\textbf{Задача 1.} В инерциальной системе отсчёта сила $F$ сообщает телу массой $m$ ускорение $a$. Во сколько раз нужно увеличить массу тела, чтобы вдвое большая сила сообщала ему в той же системе отсчёта в $4$ раза меньшее ускорение?

\kim{2}
\end{taskbasic}
\answer{8}

% === ЗАДАЧА 2 ===
% Тег: #МЕХАНИКА #КИМ_02 #ВАРИАНТ_2
\begin{taskbasic}[code={\label{task:dem26_v2_n2}}]
\textbf{Задача 2.} В инерциальной системе отсчёта сила $F$ сообщает телу массой $m$ ускорение $a$. Во сколько раз нужно уменьшить массу тела, чтобы втрое меньшая сила сообщала ему в той же системе отсчёта в $2$ раза большее ускорение?

\kim{2}
\end{taskbasic}
\answer{6}

% === ЗАДАЧА 3 ===
% Тег: #МЕХАНИКА #КИМ_02 #ВАРИАНТ_3
\begin{taskbasic}[code={\label{task:dem26_v3_n2}}]
\textbf{Задача 3.} Модуль сил гравитационного притяжения между двумя однородными шарами, центры которых находятся на расстоянии $120$~см друг от друга, равен $5$~нН. Каков будет модуль сил притяжения между ними, если расстояние между их центрами уменьшить до $30$~см?

\kim{2}
\end{taskbasic}
\answer{80}

% === ЗАДАЧА 4 ===
% Тег: #МЕХАНИКА #КИМ_02 #ВАРИАНТ_4
\begin{taskbasic}[code={\label{task:dem26_v4_n2}}]
\textbf{Задача 4.} Модуль сил гравитационного притяжения между двумя однородными шарами, центры которых находятся на расстоянии $50$~см друг от друга, равен $8$~нН. Каков будет модуль сил притяжения между ними, если расстояние между их центрами увеличить до $2$~м?

\kim{2}
\end{taskbasic}
\answer{0{,}5}

% === ЗАДАЧА 5 ===
% Тег: #МЕХАНИКА #КИМ_02 #ВАРИАНТ_5
\begin{taskbasic}[code={\label{task:dem26_v5_n2}}]
\textbf{Задача 5.} Школьник провёл исследования зависимости модуля силы трения скольжения $\vec{F}_{\text{тр}}$ от модуля нормальной составляющей силы реакции опоры $\vec{N}$. Им были получены следующие данные:
\begin{center}
\begin{tabular}{|c|c|c|c|c|}
\hline
$F_{\text{тр}}$, Н & 10 & 15 & 20 & 25 \\ \hline
$N$, Н & 50 & 75 & 100 & 125 \\ \hline
\end{tabular}
\end{center}
Определите по результатам исследования коэффициент трения скольжения.

\kim{2}
\end{taskbasic}
\answer{0{,}2}

% === ЗАДАЧА 6 ===
% Тег: #МЕХАНИКА #КИМ_02 #ВАРИАНТ_6
\begin{taskbasic}[code={\label{task:dem26_v6_n2}}]
\textbf{Задача 6.} Школьник провёл исследования зависимости модуля силы трения скольжения $F_{\text{тр}}$ от модуля нормальной составляющей силы реакции опоры $\vec{N}$. Им были получены следующие данные:
\begin{center}
\begin{tabular}{|c|c|c|c|}
\hline
$F_{\text{тр}}$, Н & 12 & 15 & 18 \\ \hline
$N$, Н & 80 & 100 & 120 \\ \hline
\end{tabular}
\end{center}
Определите по результатам исследования коэффициент трения скольжения.

\kim{2}
\end{taskbasic}
\answer{0{,}15}

% === ЗАДАЧА 7 ===
% Тег: #МЕХАНИКА #КИМ_02 #ВАРИАНТ_7
\begin{taskbasic}[code={\label{task:dem26_v7_n2}}]
\textbf{Задача 7.} Два маленьких однородных шарика массой $m$ каждый притягиваются друг к другу с гравитационными силами, равными по модулю $6$~пН. Расстояние между центрами шариков равно $r$. Каков модуль сил гравитационного притяжения друг к другу двух других маленьких однородных шариков, если масса каждого из них $3m$, а расстояние между их центрами $\frac{r}{3}$?

\kim{2}
\end{taskbasic}
\answer{54}

% === ЗАДАЧА 8 ===
% Тег: #МЕХАНИКА #КИМ_02 #ВАРИАНТ_8
\begin{taskbasic}[code={\label{task:dem26_v8_n2}}]
\textbf{Задача 8.} Два маленьких однородных шарика массой $m$ каждый притягиваются друг к другу с гравитационными силами, равными по модулю $8$~пН. Расстояние между центрами шариков равно $r$. Каков модуль сил гравитационного притяжения друг к другу двух других маленьких однородных шариков, если масса каждого из них $\frac{m}{2}$, а расстояние между их центрами $2r$?

\kim{2}
\end{taskbasic}
\answer{0{,}5}

% === ЗАДАЧА 9 ===
% Тег: #МЕХАНИКА #КИМ_02 #ВАРИАНТ_9
\begin{taskbasic}[code={\label{task:dem26_v9_n2}}]
\textbf{Задача 9.} При прямолинейном движении по горизонтальной поверхности на брусок массой $20$~кг действует сила трения скольжения, равная по модулю $20$~Н. Чему будет равен модуль силы трения скольжения после уменьшения массы тела в $2$ раза, если коэффициент трения не изменится?

\kim{2}
\end{taskbasic}
\answer{10}

% === ЗАДАЧА 10 ===
% Тег: #МЕХАНИКА #КИМ_02 #ВАРИАНТ_10
\begin{taskbasic}[code={\label{task:dem26_v10_n2}}]
\textbf{Задача 10.} При прямолинейном движении по горизонтальной поверхности на брусок массой $40$~кг действует сила трения скольжения, равная по модулю $80$~Н. Во сколько раз необходимо увеличить массу бруска, чтобы при неизменном коэффициенте трения модуль силы трения скольжения был равен $120$~Н?

\kim{2}
\end{taskbasic}
\answer{1{,}5}

% === ЗАДАЧА 11 ===
% Тег: #МЕХАНИКА #КИМ_02 #ВАРИАНТ_11
\begin{taskbasic}[code={\label{task:dem26_v11_n2}}]
\textbf{Задача 11.} В инерциальной системе отсчёта некоторая сила сообщает телу массой $5$~кг ускорение, равное по модулю $4$~$M/c^{2}$[cite: 5]. Телу какой массы вдвое большая сила сообщит в той же системе отсчёта ускорение $1$~$M/c^{2}$[cite: 5]?

\kim{2}
\end{taskbasic}
\answer{40}

% === ЗАДАЧА 12 ===
% Тег: #МЕХАНИКА #КИМ_02 #ВАРИАНТ_12
\begin{taskbasic}[code={\label{task:dem26_v12_n2}}]
\textbf{Задача 12.} В инерциальной системе отсчёта некоторая сила сообщает телу массой $2{,}5$~кг ускорение, равное по модулю $6$~$M/c^{2}$. Определите модуль ускорения, которое сообщит та же сила в той же системе отсчёта телу массой $7{,}5$~кг.

\kim{2}
\end{taskbasic}
\answer{2}

% === ЗАДАЧА 13 ===
% Тег: #МЕХАНИКА #КИМ_02 #ВАРИАНТ_13
\begin{taskbasic}[code={\label{task:dem26_v13_n2}}]
\textbf{Задача 13.} При исследовании зависимости модуля силы упругости $F_{\text{упр}}$ от удлинения пружины $\Delta x$ были получены следующие данные:
\begin{center}
\begin{tabular}{|c|c|c|c|c|}
\hline
$F_{\text{упр}}$, Н & 2,5 & 5,0 & 10,0 & 12,5 \\ \hline
$\Delta x$, м & 0,01 & 0,02 & 0,04 & 0,05 \\ \hline
\end{tabular}
\end{center}
Определите по результатам исследования жёсткость пружины.

\kim{2}
\end{taskbasic}
\answer{250}

% === ЗАДАЧА 14 ===
% Тег: #МЕХАНИКА #КИМ_02 #ВАРИАНТ_14
\begin{taskbasic}[code={\label{task:dem26_v14_n2}}]
\textbf{Задача 14.} При исследовании зависимости модуля силы упругости $F_{\text{упр}}$ от удлинения пружины $\Delta x$ были получены следующие данные:
\begin{center}
\begin{tabular}{|c|c|c|c|c|}
\hline
$F_{\text{упр}}$, Н & 1,0 & 2,0 & 3,0 & 4,0 \\ \hline
$\Delta x$, см & 0,5 & 1,0 & 1,5 & 2,0 \\ \hline
\end{tabular}
\end{center}
Определите по результатам исследования жёсткость пружины.

\kim{2}
\end{taskbasic}
\answer{200}

% === ЗАДАЧА 15 ===
% Тег: #МЕХАНИКА #КИМ_02 #ВАРИАНТ_15
\begin{taskbasic}[code={\label{task:dem26_v15_n2}}]
\textbf{Задача 15.} При исследовании зависимости модуля силы трения скольжения от модуля нормальной составляющей силы реакции опоры $\vec{N}$ были получены следующие данные:
\begin{center}
\begin{tabular}{|c|c|c|c|c|}
\hline
$F_{\text{тр}}$, Н & 0,8 & 1,6 & 2,4 & 3,2 \\ \hline
$N$, Н & 2,0 & 4,0 & 6,0 & 8,0 \\ \hline
\end{tabular}
\end{center}
Определите по результатам исследования коэффициент трения скольжения.

\kim{2}
\end{taskbasic}
\answer{0{,}4}

% === ЗАДАЧА 16 ===
% Тег: #МЕХАНИКА #КИМ_02 #ВАРИАНТ_16
\begin{taskbasic}[code={\label{task:dem26_v16_n2}}]
\textbf{Задача 16.} Ученик исследовал зависимость модуля силы трения скольжения $F_{\text{тр}}$, действующей на брусок, движущийся по горизонтальной поверхности, от его массы $m$. В таблице приведены данные, полученные в ходе эксперимента:
\begin{center}
\begin{tabular}{|c|c|c|c|c|}
\hline
$F_{\text{тр}}$, Н & 0,4 & 0,8 & 1,2 & 1,6 \\ \hline
$m$, г & 200 & 400 & 600 & 800 \\ \hline
\end{tabular}
\end{center}
Определите по результатам исследования коэффициент трения скольжения.

\kim{2}
\end{taskbasic}
\answer{0{,}2}

% === ЗАДАЧА 17 ===
% Тег: #МЕХАНИКА #КИМ_02 #ВАРИАНТ_17
\begin{taskbasic}[code={\label{task:dem26_v17_n2}}]
\textbf{Задача 17.} На рисунке изображены лабораторный динамометр и линейка. Шкала динамометра проградуирована в ньютонах, шкала линейки проградуирована в сантиметрах. Определите жёсткость пружины этого динамометра.

\kim{2}
\end{taskbasic}
\answer{10}

% === ЗАДАЧА 18 ===
% Тег: #МЕХАНИКА #КИМ_02 #ВАРИАНТ_18
\begin{taskbasic}[code={\label{task:dem26_v18_n2}}]
\textbf{Задача 18.} На рисунке изображены лабораторный динамометр и линейка. Шкала динамометра проградуирована в ньютонах, шкала линейки проградуирована в сантиметрах. Какой должна быть масса груза, подвешенного к пружине динамометра, чтобы пружина растянулась на $7{,}5$~см?

\kim{2}
\end{taskbasic}
\answer{375}

% === ЗАДАЧА 19 ===
% Тег: #МЕХАНИКА #КИМ_02 #ВАРИАНТ_19
\begin{taskbasic}[code={\label{task:dem26_v19_n2}}]
\textbf{Задача 19.} В инерциальной системе отсчёта сила величиной $35$~Н сообщает телу массой $5$~кг некоторое ускорение. Сила какой величины сообщит телу массой $9$~кг в этой же системе отсчёта такое же ускорение?

\kim{2}
\end{taskbasic}
\answer{63}

% === ЗАДАЧА 20 ===
% Тег: #МЕХАНИКА #КИМ_02 #ВАРИАНТ_20
\begin{taskbasic}[code={\label{task:dem26_v20_n2}}]
\textbf{Задача 20.} В инерциальной системе отсчёта сила величиной $70$~Н сообщает телу массой $10$~кг некоторое ускорение. Телу какой массы сила величиной $56$~Н сообщит в этой же системе отсчёта вдвое большее ускорение?

\kim{2}
\end{taskbasic}
\answer{4}